\begin{lemma}
\label{lemma-Noetherian-dvr-valuative-separation}
Let $S$ be a locally Noetherian scheme.
Let $f : X \to S$ be a morphism of schemes.
Assume $f$ is locally of finite type.
The following are equivalent:
\begin{enumerate}
\item The morphism $f$ is separated.
\item For any diagram (\ref{equation-valuative}) there is at most
one dotted arrow.
\item For all diagrams (\ref{equation-valuative}) with $A$ a discrete
valuation ring there is at most one dotted arrow.
\item For any irreducible component $X_0$ of $X$ with
generic point $\eta \in X_0$, for any discrete valuation ring
$A \subset K = \kappa(\eta)$ with fraction field $K$ and any
diagram (\ref{equation-valuative}) such that
the morphism $\Spec(K) \to X$ is the canonical one
(see Schemes, Section \ref{schemes-section-points})
there is at most one dotted arrow.
\end{enumerate}
\end{lemma}

\begin{proof}
Clearly (1) implies (2), (2) implies (3), and (3) implies (4).  It
remains to show (4) implies (1). Assume (4).
We begin by reducing to $S$ affine.  Being separated is a local
on the base (see
Schemes, Lemma \ref{schemes-lemma-characterize-separated}).
Hence, if we can show that whenever
$X \to S$ has (4) that the restriction $X_\alpha \to S_\alpha$ has (4)
where $S_\alpha \subset S$ is an (affine) open subset and $X_\alpha :=
f^{-1}(S_\alpha)$, then we will be done.  The
generic points of the irreducible components of $X_\alpha$ will be the
generic points of irreducible components of $X$, since $X_\alpha$ is
open in $X$.  Therefore, any two distinct dotted arrows in the diagram
\begin{equation}
\label{equation-valuative-alpha}
\xymatrix{
\Spec(K) \ar[r] \ar[d] & X_\alpha \ar[d] \\
\Spec(A) \ar[r] \ar@{-->}[ru] & S_\alpha
}
\end{equation}
would then give two distinct arrows in diagram
(\ref{equation-valuative}) via the maps $X_\alpha \to X$ and
$S_\alpha \to S$, which is a contradiction.  Thus we have reduced
to the case $S$ is affine.  We remark that in the course of this
reduction, we prove that if $X \to S$ has (4) then the restriction $U
\to V$ has (4) for opens $U \subset X$ and $V \subset S$ with
$f(U) \subset V$.

\medskip\noindent
We next wish to reduce to the case $X \to S$ is finite type.  Assume
that we know (4) implies (1) when $X$ is finite type. Since
$S$ is Noetherian and $X$ is locally of finite type over $S$
we see $X$ is locally Noetherian as well (see Morphisms,
Lemma \ref{morphisms-lemma-finite-type-noetherian}).
Thus, $X \to S$ is quasi-separated (see
Properties, Lemma \ref{properties-lemma-locally-Noetherian-quasi-separated}),
and therefore we may apply the valuative criterion to check whether $X$
is separated (see
Schemes, Lemma \ref{schemes-lemma-valuative-criterion-separatedness}).
Let $X = \bigcup_\alpha X_\alpha$ be an affine open
cover of $X$. Given any two dotted arrows, in a diagram
(\ref{equation-valuative}), the image of the closed points of
$\Spec A$ will
fall in two sets $X_\alpha$ and $X_\beta$.  Since $X_\alpha \cup
X_\beta$ is open, for topological reasons it must contain the image of
$\Spec(A)$ under both maps. Therefore, the two dotted arrows factor
through $X_\alpha \cup X_\beta \to X$, which is a scheme of finite type over
$S$. Since $X_\alpha \cup X_\beta$ is an open subset of $X$, by our
previous remark, $X_\alpha \cup X_\beta$ satisfies (4), so by
assumption, is separated.  This implies the two given dotted
arrows are the same. Therefore, we have reduced to $X \to S$ is finite type.

\medskip\noindent
Assume $X \to S$ of finite type and assume (4).
Since $X \to S$ is finite type, and $S$ is an affine Noetherian
scheme, $X$ is also Noetherian (see
Morphisms, Lemma \ref{morphisms-lemma-finite-type-noetherian}).
Therefore, $X \to X \times_S X$ will
be a quasi-compact immersion of Noetherian schemes.  We proceed by
contradiction.  Assume that $X \to X \times_S X$ is not closed.  Then,
there is some $y \in X \times_S X$ in the closure of the image that is
not in the image. As $X$ is Noetherian it has finitely many irreducible
components. Therefore, $y$ is in the closure of the image of one of
the irreducible components $X_0 \subset X$.  Give $X_0$ the reduced
induced structure.  The composition $X_0 \to X \to X \times_S X$
factors through the closed subscheme $X_0 \times_S X_0 \subset X \times_S X$.
Denote the closure of $\Delta(X_0)$ in $X_0 \times_S X_0$
by $\bar X_0$ (again as a reduced closed subscheme). Thus $y \in \bar X_0$.
Since $X_0 \to X_0 \times_S X_0$ is an immersion, the image of $X_0$
will be open in $\bar X_0$. Hence $X_0$ and $\bar X_0$ are
birational. Since $\bar{X}_0$ is a closed subscheme of a
Noetherian scheme, it is Noetherian. Thus, the local ring
$\mathcal O_{{\bar X_0, y}}$ is a local Noetherian domain with fraction
field $K$ equal to the function field of $X_0$.  By the Krull-Akizuki
theorem (see Algebra, Lemma \ref{algebra-lemma-exists-dvr}), there exists a
discrete valuation ring $A$ dominating $\mathcal O_{{\bar X_0, y}}$
with fraction field $K$.  This allows to construct a diagram:
\begin{equation}
\label{equation-valuative-generic}
\xymatrix{
\Spec(K) \ar[r] \ar[d] & X_0 \ar[d]^{\Delta} \\
\Spec(A) \ar[r] \ar@{-->}[ur]& X_0 \times_S X_0 \\
}
\end{equation}
which sends $\Spec K$ to the generic point of $\Delta(X_0)$ and
the closed point of $A$ to $y \in X_0 \times_S X_0$ (use the material in
Schemes, Section \ref{schemes-section-points} to construct the arrows).
There cannot even exist
a set theoretic dotted arrow, since $y$ is not in the image of
$\Delta$ by our choice of $y$.  By categorical means, the existence of
the dotted arrow in the above diagram is equivalent to the uniqueness
of the dotted arrow in the following diagram:
\begin{equation}
\label{equation-valuative-nonexistent}
\xymatrix{
\Spec(K) \ar[r] \ar[d] & X_0 \ar[d]\\
\Spec(A) \ar[r] \ar@{-->}[ur] & S \\
}
\end{equation}
Therefore, we have non-uniqueness in this latter diagram by the
nonexistence in the first.  Therefore, $X_0$ does not satisfy
uniqueness for discrete valuation rings, and since $X_0$ is an
irreducible component of $X$, we have that $X \to S$ does not satisfy
(4).  Therefore, we have shown (4) implies (1).
\end{proof}
